% ---
% title: KL Divergence in Reinforcement Learning
% date: 2026/2/13
% tag: RL
% ---
\section{不同的KL散度}

There are two types of KL divergence in Reinforcement Learning, \textbf{forward KL} and \textbf{reverse KL}. In this post, we will focus on the \textbf{Loss Difference} between forward KL and reverse KL.

\subsection{\textbf{一、KL 散度}不是对称的}

\[
\mathrm{KL}(p \,\|\, q) \;\neq\; \mathrm{KL}(q \,\|\, p)
\]

在 LLM-RL 里有两个自然的选择：

\subsubsection{Forward KL}

\[
\mathrm{KL}(\pi_{\text{base}} \,\|\, \pi_\theta)
\]

\subsubsection{Reverse KL}

\[
\mathrm{KL}(\pi_\theta \,\|\, \pi_{\text{base}})
\]

\subsection{二、为什么 RL 里用的是 reverse KL？}

引用的这句话点中了核心：

\begin{quote}
\emph{so that the learned policy assigns high probability mass to a narrow set of high-reward trajectories}
\end{quote}

这正是 \textbf{reverse KL 的典型性质}。

\subsubsection{🔹 Reverse KL 的“模式寻优（mode-seeking）”行为}

当你最小化：

\[
\mathrm{KL}(\pi_\theta \,\|\, \pi_{\text{base}}) = \mathbb{E}_{\pi_\theta}\left[\log \frac{\pi_\theta}{\pi_{\text{base}}}\right]
\]

意味着：

\begin{itemize}
  \item 期望是 \textbf{在当前策略} \(\pi_\theta\) \textbf{下}
  \item 只关心 \(\pi_\theta\) 已经采样到的区域
  \item \textbf{允许忽略 base policy 覆盖但 reward 低的区域}
\end{itemize}

结果就是：

\begin{itemize}
  \item 把概率集中到\textbf{少量高 reward 序列}
  \item 快速“塌缩”到 preferred behaviors
  \item 非常适合 RL 的目标（maximize expected reward）
\end{itemize}

\subsection{三、为什么不用 forward KL？}

对比一下 forward KL：

\[
\mathrm{KL}(\pi_{\text{base}} \,\|\, \pi_\theta) = \mathbb{E}_{\pi_{\text{base}}}\left[\log \frac{\pi_{\text{base}}}{\pi_\theta}\right]
\]

\subsubsection{🔹 Forward KL 的“覆盖优先（mode-covering）”行为}

它会强迫：

\begin{itemize}
  \item \(\pi_\theta\) 在 \textbf{base policy 有概率的所有地方都给概率}
  \item 即使这些地方 reward 很低
\end{itemize}

结果是：

\begin{itemize}
  \item 保守
  \item reward 提升慢
  \item \textbf{牺牲性能换取分布覆盖}
\end{itemize}

\hrule

论文里说的 \textbf{reverse KL} 是：

\[
\mathrm{KL}(\pi_\theta \,\|\, \pi_{\text{base}}) \;=\; \mathbb{E}_{x \sim \pi_\theta} \left[ \log \frac{\pi_\theta(x)}{\pi_{\text{base}}(x)} \right]
\]

这里的 x 是：\textbf{一个完整序列}（whole trajectory / whole response）

不是 token，也不是 prefix。

\subsection{把期望写成“显式求和”}

\[
\mathrm{KL}(\pi_\theta \,\|\, \pi_{\text{base}}) = \sum_{x \in \mathcal{X}} \pi_\theta(x) \left( \log \pi_\theta(x) - \log \pi_{\text{base}}(x) \right)
\]

\subsubsection{⚠️ 注意：这里的求和对象}

\begin{itemize}
  \item \(\mathcal{X}\)：\textbf{所有可能的 token 序列}
  \item 长度可变
  \item 每一步 vocab size \textasciitilde{} 30k–100k
\end{itemize}

\hrule

\section{計算Reverse KL 散度的一般方法**}

\subsection{一、問題}

\textbf{1️⃣ 序列空间是指数级的}

假设：vocab size = V = \(50\,000\), 最大长度 = T = \(512\)。那么序列空间大小：

\[
\mathcal{X}| = \sum_{t=1}^{512} V^t \;\approx\; V^{512}
\]

这是一个\textbf{天文数量}，\textbf{不可能枚举所有序列，}也不可能对它们做显式加权求和

\textbf{2️⃣ 和 token-level logprob 的“可计算性”形成强烈对比}

给定一个具体序列\(x = (a_1,\dots,a_T)\)：

\[
\log \pi_\theta(x) = \sum_{t=1}^T \log \pi_\theta(a_t \mid a_{<t})
\]

这一步是：

\begin{itemize}
  \item 前向传播
  \item teacher forcing
  \item \textbf{完全可算}
\end{itemize}

同理\(\log \pi_{\text{base}}(x)\)也是可算的。算的是 \textbf{“给定一个序列，算它的 log-prob”，而}不可算的是 \textbf{“对所有序列取期望”}

\subsection{二、为什么 forward KL 反而“更像可算的”？}

对比 forward KL：

\[
\mathrm{KL}(\pi_{\text{base}} \,\|\, \pi_\theta) = \mathbb{E}_{x \sim \pi_{\text{base}}} \left[ \log \frac{\pi_{\text{base}}(x)}{\pi_\theta(x)} \right]
\]

关键差别在采样分布**

\begin{itemize}
  \item \(\pi_{\text{base}}\)：\textbf{固定的}
  \item 可以离线采样
  \item 可以用数据集近似
\end{itemize}

但：

\begin{itemize}
  \item \(\pi_\theta\)：\textbf{训练中不断变化}
  \item 每次更新分布都变
  \item 无法预先覆盖 support
\end{itemize}

📌 这也是为什么 reverse KL 更“RL 风格”，但更难处理。

\subsection{三、那 PPO / RLHF 里到底是怎么算的？}

现实中大家做的是：

\[
\mathrm{KL}(\pi_\theta \,\|\, \pi_{\text{base}}) \;\approx\; \frac{1}{N} \sum_{i=1}^N \left( \log \pi_\theta(x^{(i)}) - \log \pi_{\text{base}}(x^{(i)}) \right), \quad x^{(i)} \sim \pi_\theta
\]

也就是：

\begin{itemize}
  \item 从当前 policy 采样 \(N\) \textbf{个序列}
  \item 用 \(\color{red}Monte\ Carlo\) 估计期望
\end{itemize}

\[
\frac{1}{N}\sum_{i=1}^N f(x^{(i)}),
\quad x^{(i)} \sim \pi_\theta
\;\;\xrightarrow[N\to\infty]{}\;\;
\mathbb{E}_{x\sim\pi_\theta}[f(x)]
\]

📌 这就是 \emph{sample-based estimators of the reverse KL divergence}

\subsection{四、但注意：“数值可估 ≠ 梯度无偏”}

这是你现在这篇论文最重要的洞察：

\begin{itemize}
  \item 即便你能 MC 估计\(\mathrm{KL} \approx \hat{K}\)
  \item \(\nabla_\theta \hat{K}≠\nabla_\theta \mathrm{KL}\)
\end{itemize}

除非 estimator 的结构是对的。

这也是为什么：

\begin{itemize}
  \item 放在 \textbf{reward} 里 → OK（score function trick）
  \item 放在 \textbf{loss} 里 → 梯度错位
\end{itemize}
